\section{Study Assessments and Procedures}

This section specifies the efficacy and safety assessments, the device-telemetry
measurements unique to the Physical AI component, and the adverse-event and
unanticipated-problem machinery for the study. The timing of every assessment is
governed by the Schedule of Activities ($\S$1.3); this section defines what is
measured, how it is graded, and how it is reported. Because the study is conducted
across eight academic high-volume hepatopancreatobiliary centers, every endpoint
that admits central reading is read centrally: a central imaging core performs the
blinded independent central review (BICR), a central pathology core performs the
masked synoptic margin and regression reads, and central laboratories perform the
biomarker and clinical-chemistry assays, so that ascertainment is uniform across
sites and independent of the open-label assignment.

\subsection{Efficacy Assessments}

The primary efficacy endpoint is progression-free survival (PFS), the time from
randomization to the first documented radiographic disease progression by RECIST
version 1.1 or death from any cause, whichever occurs first. Imaging is acquired on
contrast-enhanced cross-sectional studies at the protocol-specified timepoints and is
read by the central imaging core; the investigator-assessed read is the primary
analysis and the BICR read supports a prespecified sensitivity analysis ($\S$9.4.2).
Overall survival (OS), the time from randomization to death from any cause, is
followed to the 24-month endpoint and is a key secondary endpoint, with the 2025
Dutch nationwide robotic pancreaticoduodenectomy cohort retained as an external
descriptive reference \cite{DutchCohort2025}.

The surgical and biomarker efficacy endpoints are assessed centrally. The
margin-negative (R0) resection rate is determined by central masked synoptic
pathology, with the pathologist masked to the assigned arm and to the device
telemetry and with masked independent adjudication of discordant or indeterminate
margins. The major pathologic response (MPR) is graded by the same central pathology
core on the resection specimen. Postoperative pancreatic fistula is graded by the
International Study Group of Pancreatic Surgery (ISGPS) definition as a biochemical
leak (no longer classified as a clinically relevant fistula), Grade B (requiring a
change in management), or Grade C (requiring reoperation or causing organ failure or
death) \cite{Bassi2017}; the clinically relevant (Grade B or C) rate is a key
secondary endpoint. Circulating tumor DNA (ctDNA) bearing the participant's KRAS
mutation is measured by a central laboratory, and week-12 molecular clearance, the
conversion from a detectable to an undetectable KRAS ctDNA level by week 12, is a key
secondary endpoint \cite{Tie2019ctdna}.

Anastomosis quality is a prespecified Arm A device-quality assessment and is recorded
for each of the three reconstructions against its controller ring-tension band: the
duct-to-mucosa pancreaticojejunostomy (0.30 to 0.60 N), the end-to-side
hepaticojejunostomy (0.20 to 0.50 N), and the antecolic gastrojejunostomy (0.40 to
0.80 N), as constructed under continuous Arm A telemetry ($\S$6). The
pancreaticojejunostomy ring tension is the dominant fistula determinant and is
recorded against the fistula and biomarker reads above. Health-related quality of
life is measured with the EORTC QLQ-C30 core questionnaire and the PAN26
pancreatic-cancer module at the protocol-specified timepoints \cite{EORTCqlqc30}.

\subsection{Safety and Other Assessments}

Clinical safety assessments comprise vital signs, a directed physical examination,
hematology, serum chemistry including hepatic and pancreatic indices, coagulation
studies, the serum daraxonrasib trough assay for Arm A, and protocol-specified
imaging, at the frequencies in the Schedule of Activities, with clinical-chemistry
and biomarker assays routed to the central laboratories. In addition to these
conventional assessments, the Arm A device generates a continuous safety telemetry
stream that is a source of protocol observations under 21 CFR $\S$312.23(g): per-arm
and cumulative force profiles against the $\leq 3$ N per-arm and $\leq 18$ N
cumulative cross-arm caps, trajectory accuracy against specification, every vascular
safety-zone verdict, every heartbeat and watchdog event, and every E-stop
activation, each captured at its recorded rate in the hash-chained audit trail.
Pre-procedure safety-matrix results are recorded for each Arm A case.

The vascular safety-zone gate ticks at 10 kHz and evaluates the instrument-tip
proximity to each of five named vessels against three concentric thresholds, with
the verdict driving the response: a soft-warning zone at 3.0 mm raises an LLM
advisory, a no-fly zone (1.0 mm at the superior mesenteric and portal veins; 1.5 mm
at the hepatic artery, celiac axis, and superior mesenteric artery) blocks the
command, and a hard-stop zone at 5.0 mm forces a $\leq 3$ ms cross-arm E-stop. The
gate exercises all five vessels across all eight phases, and every
verdict is written to the audit trail. Gate is shown in Figure 15.

\begin{mermaidfig}
\node[mmin]  (tip)  at (4.7,0)     {Instrument-tip\\proximity\\sampled at 10 kHz};
\node[mmdec] (z)    at (4.7,-2.4)  {Zone test\\per vessel};
\node[mmin]  (clr)  at (0,-4.8)    {Clear\\continue command};
\node[mmstep](soft) at (4.7,-4.8)  {soft\_warning\\LLM advisory\\all vessels 3.0 mm};
\node[mmstep](nofly)at (9.6,-4.8)  {no\_fly: block\\SMV/PV 1.0 mm\\HA/CA/SMA 1.5 mm};
\node[mmgoal](hard) at (9.6,-2.4)  {hard\_stop: ESTOP\\all vessels 5.0 mm\\halt $\leq 3$ ms};
\node[mmdark](v)    at (0,-2.4)    {Vessels: SMV, PV,\\hepatic artery,\\celiac axis, SMA\\(5 $\times$ 8 phases)};
\draw[mmedge] (tip) -- (z);
% \draw[mmedge] (z) -- node[left,font=\scriptsize, xshift=-0.2cm]{clear} (clr);
% \draw[mmedge] (z) -- node[right,font=\scriptsize,pos=0.6, yshift=0.2cm]{soft} (soft);
% \draw[mmedge] (z) -- node[right,font=\scriptsize,pos=0.6, xshift=0.2cm]{no\_fly} (nofly);
% \draw[mmedge] (z) -- node[above,font=\scriptsize, xshift=-0.05cm]{hard} (hard);

\draw[mmedge] (z.south west) -- node[left,font=\scriptsize, xshift=-0.2cm]{clear} (clr.north east);
\draw[mmedge] (z) -- node[right,font=\scriptsize,pos=0.6, yshift=0.2cm]{soft} (soft);
\draw[mmedge] (z.south east) -- node[right,font=\scriptsize,pos=0.6, xshift=0.2cm]{no\_fly} (nofly.north west);
\draw[mmedge] (z) -- node[above,font=\scriptsize, xshift=-0.05cm]{hard} (hard);

\draw[mmedged] (z) -- (v);
\end{mermaidfig}
\vspace{-0.7cm}
\figcaption{Figure 15. Five-vessel vascular safety-zone gate for Arm A. Three
concentric proximity zones around each\\vessel drive response: a
soft-warning zone raises an LLM advisory, a no-fly zone (superior mesenteric and\\portal veins; hepatic artery, celiac axis, and
superior mesenteric artery) blocks the command, and a hard\\stop zone forces a cross-arm E-stop. The gate is sampled at 10 kHz across all eight
phases of the operation.}

The halt chain that backs the gate is layered. The 10 kHz heartbeat bus broadcasts
a 64-byte frame on a $100\,\mu$s per-arm watchdog deadline; a frame that is on time
continues the command state, whereas a missed or out-of-parameter frame parks the
affected arm within $50\,\mu$s and escalates to a cross-arm E-stop that halts the
platform in $\leq 3$ ms at a 0.0 N force clamp. A software-independent hardware
E-stop meeting the $\leq 500$ ms requirement backs the software chain. The
architecture is shown in Figure 16.

\subsection{Adverse Events and Serious Adverse Events}

Adverse events are captured through two parallel and equally binding streams: a
conventional pharmacovigilance stream for clinical events in both arms, and a
Physical AI safety stream for device- and AI-related events in Arm A. The combined
workflow is shown in Figure 17.

\clearpage

\begin{mermaidfig}
\node[mmgoal](hb)  at (4.7,0)     {10 kHz heartbeat bus\\64-byte frame,\\$100\,\mu$s deadline};
\node[mmdec] (wd)  at (4.7,-2.4)  {Watchdog window\\$100\,\mu$s per arm};
\node[mmin]  (ok)  at (0,-2.4)    {Frame on time\\continue command\\state};
\node[mmstep](miss)at (9.4,-2.4)  {Frame missed or\\out-of-parameter};
\node[mmstep](park)at (9.4,-4.7)  {Emergency arm park\\within $50\,\mu$s};
\node[mmgoal](est) at (4.7,-4.7)  {Cross-arm ESTOP\\halt $\leq 3$ ms,\\force clamp 0.0 N};
\node[mmdark](hw)  at (0,-4.7)    {Hardware E-stop\\independent of\\software ($\leq 500$ ms)};
\draw[mmedge] (hb) -- (wd);
\draw[mmedge] (wd) -- node[above,font=\scriptsize, xshift=0.1cm]{met} (ok);
\draw[mmedge] (wd) -- node[above,font=\scriptsize, xshift=-0.1cm]{missed} (miss);
\draw[mmedge] (miss) -- (park);
\draw[mmedge] (park) -- (est);
\draw[mmedge] (est) -- (hw);
\end{mermaidfig}
\vspace{-0.65cm}
\figcaption{Figure 16. Heartbeat, watchdog, and E-stop architecture for Arm A. A
10 kHz heartbeat bus with a $100\,\mu$s\\per arm watchdog continues the command state
on an on-time frame; a missed frame parks the arm\\within $50\,\mu$s and escalates to
a cross-arm E-stop that halts the platform in $\leq 3$ ms at a 0.0 N force\\clamp, backed by a software-independent hardware E-stop meeting the $\leq 500$ ms
requirement.}

\begin{mermaidfig}
\node[mmin]  (det)  at (0,0)     {Event detected\\clinical AE or\\Physical AI AE};
\node[mmdec] (cls)  at (5.2,-2.4)  {Classify severity /\\causality /\\expectedness};
\node[mmstep](ns)   at (0,-2.4)    {Non-serious\\record in CRF\\(CTCAE v5)};
\node[mmgoal](s15)  at (10.4,-2.4)  {Serious +\\unexpected\\report $\leq 15$ days};
\node[mmgoal](s7)   at (10.4,0)     {Fatal / life-\\threatening\\report $\leq 7$ days};
\node[mmstep](pai)  at (0,-4.8)    {Physical AI triggers\\system-drug, cyber,\\degradation, $> 2\times$};
\node[mmdark](aud)  at (5.2,-4.8)  {Preserve audit trail\\$-24$ h to $+72$ h\\federated, all sites};
\node[mmgoal](comm) at (10.4,-4.8)  {Notify FDA, IRB,\\DSMB, Safety\\Review Committee};
\draw[mmedge] (det.east) -- ++(1.26,0) |- (cls.north);
\draw[mmedge] (cls) -- node[above,font=\scriptsize, xshift=0.15cm, yshift=-0.7cm]{non-\\serious} (ns);
\draw[mmedge] (cls) -- node[above,font=\scriptsize,pos=0.6, xshift=-0.22cm, yshift=-0.4cm]{serious} (s15);
\draw[mmedge] (cls.east) -- ++(0.5,0) |- (s7.west);
\draw[mmedge] (det.west) -- ++(-0.5,0) |- (pai.west);
\draw[mmedge] (pai) -- (aud);
\draw[mmedge] (aud) -- (comm);
\draw[mmedge] (s15) -- (comm);
\end{mermaidfig}
\vspace{-0.65cm}
\figcaption{Figure 17. Adverse-event and Physical AI AE reporting workflow. Clinical
events are graded by CTCAE v5 and reported on the 15-day (serious, unexpected) or
7-day (fatal, life-threatening) IND/IDE timelines, while Physical AI triggers add
their own reports, each preserving the hash-chained audit trail from 24 hours before
to 72 hours after the event, federated across sites, with notification to FDA, IRB, DSMB, and Physical AI Safety Review Committee.}

% \clearpage

\subsubsection{Definition of Adverse Events (AE)}

An adverse event is any untoward medical occurrence in a participant administered the
study intervention, whether or not considered related to it; the term encompasses
any unfavorable and unintended sign (including an abnormal laboratory finding),
symptom, or disease temporally associated with the intervention. Clinical adverse
events in both arms are graded for severity by the Common Terminology Criteria for
Adverse Events (CTCAE) version 5.0. In addition, a Physical AI adverse event, per 21
CFR $\S$312.32(f), is any untoward occurrence associated with the operation of the
Arm A device during the investigation, whether or not caused by it, including robotic
malfunction, AI prediction error, sensor failure, communication loss, unauthorized
system access, digital-twin divergence, and any event requiring activation of the
emergency stop \cite{kawchak2026adapt312}.

\subsubsection{Definition of Serious Adverse Events (SAE)}

A serious adverse event is any adverse event that results in death, is
life-threatening, requires inpatient hospitalization or prolongation of existing
hospitalization, results in persistent or significant disability or incapacity, is a
congenital anomaly or birth defect, or is an important medical event that may
jeopardize the participant and may require intervention to prevent one of the
foregoing outcomes. A serious Physical AI adverse event, per $\S$312.32(f), is a
Physical AI adverse event that results in, or has the potential to result in, any
serious-adverse-event outcome, or that results in an unplanned interruption of an
investigational drug-administration procedure requiring clinical intervention, an
uncontrolled release of the investigational drug, or compromise of the device's
sterile field during a surgical procedure.

% \clearpage

\subsubsection{Classification of an Adverse Event}

Each adverse event is classified along four axes. Severity is graded by CTCAE
version 5.0 (Grade 1 mild through Grade 5 fatal). Causality is assessed separately
in relation to the drug and in relation to the device, each as related or unrelated
(with relatedness denoting a reasonable possibility that the component caused the
event), so that a single event may carry independent drug and device causality
attributions; in Arm B, which has no investigational device, causality is assessed
in relation to the drug and to the procedure. Expectedness is judged against the
reference safety information (the daraxonrasib investigator's brochure for the drug
and the Physical AI supplement and System Description for the device), and an event
not listed there at the observed specificity or severity is classified as
unexpected. Seriousness is determined against the criteria in $\S$8.3.2. The
combination of unexpected, serious, and at least possibly related defines a suspected
adverse reaction reportable on the expedited timelines below.

\subsubsection{Time Period and Frequency for Assessment and Follow-Up}

Adverse events are solicited and recorded from the time of informed consent and are
assessed at every study contact. The primary intensive safety-collection window
spans from consent through the 30-day postoperative safety visit. Beyond the 30-day
window, serious adverse events, events of special interest, Physical AI adverse
events, and new anticancer therapy continue to be collected through the long-term
follow-up period to the 24-month overall-survival endpoint. Each event is followed
until resolution, stabilization, return to baseline, or a determination that no
further change is expected, or until the participant is lost to follow-up or
withdraws consent.

\subsubsection{Adverse Event Reporting}

All adverse events, serious and non-serious, are recorded in the case report form
with onset and resolution dates, severity grade, drug and device causality,
expectedness, seriousness, action taken with each intervention, and outcome. Routine
(non-serious) events are captured on the adverse-event case report form at the
scheduled visits and do not require expedited transmission. Device-telemetry-derived
observations in Arm A that meet the adverse-event definition are entered with
reference to the corresponding audit-trail records.

\clearpage

\subsubsection{Serious Adverse Event Reporting}

Serious adverse events are reported by the investigator to the sponsor within
24 hours of awareness. The sponsor submits IND/IDE safety reports to the FDA on the
regulatory timelines: any unexpected fatal or life-threatening suspected adverse
reaction is reported as soon as possible but no later than 7 calendar days after
initial receipt of the information, and any other serious and unexpected suspected
adverse reaction is reported no later than 15 calendar days after the determination
that the event is reportable. In addition to the clinical SAE pathway, six Physical
AI reporting triggers are reported under $\S$312.32(g): (i) any serious Physical AI
adverse event, on the 7- or 15-day timeline as applicable; (ii) any Physical AI
system-drug interaction event causing incorrect drug administration (wrong dose,
route, timing, or patient), regardless of seriousness; (iii) any cybersecurity
incident, within 15 calendar days (or 7 days if it has the potential to cause patient
harm); (iv) sustained AI model performance degradation below the pre-specified
acceptance criteria for three or more consecutive procedures or a cumulative 24
hours; (v) any sim-to-real divergence beyond twice the validated gap tolerance
(trajectory deviation $> 3$ mm or force discrepancy $> 0.8$ N); and (vi) any
digital-twin discrepancy beyond the pre-specified accuracy thresholds. For every
Physical AI adverse event reported, the complete hash-chained audit trail covering
the period from 24 hours before the event to 72 hours after the event is preserved
under 21 CFR part 11, federated across the eight sites, and made available to the FDA
on request \cite{ecfr2026part11}.

\subsubsection{Events of Special Interest}

The following events are designated events of special interest and are tracked and
reported with the same diligence as serious adverse events regardless of whether they
independently meet seriousness criteria: vascular injury (including any hard-stop
E-stop precipitated by a vascular safety-zone breach), clinically relevant
postoperative pancreatic fistula (ISGPS Grade B or C) \cite{Bassi2017}, conversion
from device-directed to manual or open technique, any cybersecurity incident
affecting the device, and any model-performance degradation. Each is reviewed by the
Physical AI Safety Review Committee, which convenes at intervals not exceeding 90 days
during active enrollment, and is reported to the data safety monitoring board (DSMB).

\subsubsection{Reporting of Pregnancy}

Pregnancy is not an adverse event but is reported because of the potential exposure
of an embryo or fetus to the study chemotherapy. Participants of childbearing
potential and partners of participants must use highly effective contraception during
treatment and for the protocol-specified interval afterward. Any pregnancy in a
participant or a participant's partner occurring during treatment or within that
interval is reported by the investigator to the sponsor within 24 hours of awareness
and is followed to outcome; a pregnancy with an associated serious adverse event, or
an abnormal pregnancy outcome, is reported on the applicable expedited SAE timeline.

\clearpage

\subsection{Unanticipated Problems}

\subsubsection{Definition}

An unanticipated problem is any incident, experience, or outcome that is unexpected
(not described in the protocol, consent, investigator's brochure, or Physical AI
System Description in terms of nature, severity, or frequency given the procedures
and population), related or possibly related to participation, and that suggests the
research places participants or others at a greater risk of harm (including physical,
psychological, economic, or social harm) than was previously known or recognized. For
the device, the Physical AI hold grounds (repeated uncorrected safety failures, a
confirmed cybersecurity compromise of drug-administration parameters or system
integrity, and a simulation-validation deficiency that undermines the pre-clinical
safety basis) are treated as paradigm unanticipated problems.

\subsubsection{Reporting}

Unanticipated problems are reported by the investigator to the reviewing IRB and, as
applicable, to the sponsor and the FDA in accordance with the institutional policy
and the IND/IDE requirements. An unanticipated problem that is also a serious and
unexpected suspected adverse reaction is reported on the 7- or 15-day timeline of
$\S$8.3.6; other unanticipated problems are reported promptly and ordinarily within
the institution's defined window (typically no later than the next scheduled
continuing review, and sooner when the problem indicates an immediate hazard).
Reports that involve the device preserve the hash-chained audit trail across the
event window, federated across the eight sites \cite{ecfr2026part11}. The Physical AI
Safety Review Committee and the DSMB review each unanticipated problem and recommend
any change to device operation, human-oversight requirements, Unified Safety Level
thresholds, or the protocol.

\subsubsection{Reporting to Participants}

When an unanticipated problem affects the rights, safety, or welfare of enrolled
participants, or could affect a participant's willingness to continue, the
information is communicated to affected participants in a timely manner, ordinarily
through an IRB-approved re-consent or written notification. Where an unanticipated
problem changes the risk profile of the combined intervention, including a relevant
Physical AI safety event, the consent document is updated and currently enrolled
participants are re-consented before any further study procedure, consistent with the
Physical AI opt-out in the original consent.
